\section[A degree-27 chart presentation]{A degree-\(27\) chart presentation}
\label{sec:chart-scheme}

\subsection[The U01 equations]{The \(U_{01}\) equations}

On the standard Pl\"ucker chart \(U_{01}\), write a line as the row span of
\begin{equation}
 M(a,b,c,d)=
 \begin{pmatrix}
 1&0&a&b\\
 0&1&c&d
 \end{pmatrix}.
\label{eq:chart-matrix}
\end{equation}
With parameters \(s,t\), its points are
\[
 (s,t,as+ct,bs+dt).
\]
Restricting \cref{eq:cubic} gives
\begin{equation}
 F(s,t,as+ct,bs+dt)
 =f_0s^3+f_1s^2t+f_2st^2+f_3t^3.
\label{eq:restricted-cubic}
\end{equation}
The four sparse polynomials \(f_i\in\Q[a,b,c,d]\) are printed in
\cref{app:chart-data}.  Their common zero scheme is
\(\FanoY\cap U_{01}\).

The degree-order calculation has a \(21\)-element Gr\"obner basis.  Its
standard monomials have total-degree counts
\begin{equation}
  (H_0,H_1,H_2,H_3,H_{\ge4})=(1,4,10,12,0),
\label{eq:chart-hilbert}
\end{equation}
so the quotient dimension is \(27\).  Exact FGLM conversion gives a
four-element lexicographic shape
\begin{equation}
\begin{aligned}
 g(d)&=0,\\
 \lambda_c c+h_c(d)&=0,\\
 \lambda_b b+h_b(d)&=0,\\
 \lambda_a a+h_a(d)&=0,
\end{aligned}
\qquad
\lambda_a\lambda_b\lambda_c\ne0,
\label{eq:lex-shape}
\end{equation}
where \(g\in\mathbf Z[d]\) is primitive of degree \(27\), and each
\(h_\bullet\) has degree at most \(26\).  The complete coefficient arrays
are carried by the exact electronic certificate; the compact invariants and
the coefficient endpoints of \(g\) are recorded in
\cref{app:chart-data}.

The shape computation alone is not used as an ideal-membership proof.
Instead, the independent replay substitutes
\begin{equation}
 a=-\frac{h_a(d)}{\lambda_a},\qquad
 b=-\frac{h_b(d)}{\lambda_b},\qquad
 c=-\frac{h_c(d)}{\lambda_c}
\label{eq:back-substitution}
\end{equation}
into each \(f_i\), clears the three constant denominators, and reduces the
resulting numerator modulo \(g\).  All four remainders are exactly zero.

\subsection{From the chart to the global scheme}

Let
\[
 A_g=\Q[d]/(g).
\]
The four direct remainder identities induce a surjection
\begin{equation}
 \Q[a,b,c,d]/(f_0,f_1,f_2,f_3)
 \twoheadrightarrow A_g,
\label{eq:chart-surjection}
\end{equation}
because the image contains the class of \(d\).  Dually, this gives a closed
immersion
\begin{equation}
 Z:=\Spec(A_g)\hookrightarrow \FanoY\cap U_{01}.
\label{eq:chart-closed}
\end{equation}

\begin{lemma}[Open-and-closed bridge]
\label{lem:clopen}
The immersion in \cref{eq:chart-closed} is a closed immersion into the global
scheme \(\FanoY\).
\end{lemma}

\begin{proof}
By \cref{prop:fano-etale}, \(\FanoY\) is a finite \(\etale\) scheme over a
field.  Its underlying space is finite and discrete, and every open
subscheme is a union of connected components.  Thus
\(\FanoY\cap U_{01}\) is open-and-closed in \(\FanoY\).  Composing
\cref{eq:chart-closed} with this closed inclusion proves the claim.
\end{proof}

\begin{proposition}[Complete chart presentation]
\label{prop:complete-chart}
There is an isomorphism
\[
  \FanoY\cong\Spec\!\bigl(\Q[d]/(g)\bigr).
\]
In particular all \(27\) geometric lines lie in \(U_{01}\).
\end{proposition}

\begin{proof}
The algebra \(A_g\) has \(\Q\)-dimension \(27\), because
\(\deg g=27\).  Hence \(Z\) and \(\FanoY\) are finite schemes of the same
rank.  The global closed immersion from \cref{lem:clopen} corresponds to a
surjective map between their \(27\)-dimensional coordinate algebras.  It is
therefore an isomorphism.
\end{proof}

\subsection{Independent complement guard}

The exact replay also generates the other five standard charts and adds the
equation \(p_{01}=0\), which cuts out the complement of \(U_{01}\).  Each
resulting ideal has reduced Gr\"obner basis \(\{1\}\):
\begin{center}
\begin{tabular}{@{}ccc@{}}
\toprule
chart & row-coordinate convention & equation for \(p_{01}=0\)\\
\midrule
\(U_{02}\)&\((s,as+ct,t,bs+dt)\)&\(c=0\)\\
\(U_{03}\)&\((s,as+ct,bs+dt,t)\)&\(c=0\)\\
\(U_{12}\)&\((as+ct,s,t,bs+dt)\)&\(c=0\)\\
\(U_{13}\)&\((as+ct,s,bs+dt,t)\)&\(c=0\)\\
\(U_{23}\)&\((as+ct,bs+dt,s,t)\)&\(ad-bc=0\)\\
\bottomrule
\end{tabular}
\end{center}
This calculation independently audits the chart convention.  It is not used
circularly to justify \cref{lem:clopen}; global equality already follows
from finite \(\etale\) rank comparison.
